\section{Discussion and Conclusion}
The above results provide support for the proposition that at least some of the value of cultural variety in contemporary American society is \emph{social}: the ability to appropriate a wide variety of symbolic goods allows the individual to maintain---and diachronically replenish \citep{cardon_granjon05, lazarsfeld_merton54}---\emph{a larger} and \emph{a more varied} personal network.   Persons who engage in a variety of cultural activities are more likely to have successfully mobilized a socially embedded strategy when looking for a job, in particular those involving the activation of \emph{weak ties} but not those involving the activation of strong or embedded ties.  Most importantly this tie mobilization appears to reflect precisely the type of network advantages that are associated with cultural variety: the activation of weaker arms-length relationships. These are precisely those relationships which in certain contexts---such as the loosely-knit networks of constantly shifting relationships associated typical of high-status Euro-American Westerners---that are most likely to produce advantages in the labor market (by for instance providing information about high-paying jobs or jobs that provide more of a fit between personal competences and qualifications and actual task activities). 

In addition, I find that while engaging in a wide variety of activities increases the chances of engaging not all cultural activities have the same functional linkage to network activation. Activities closer to the category of ``popular culture'' (e.g. mass disseminated via ``culture-industry systems'') appears to increase the chances of having successfully learned about the current job via an acquaintance; this is consistent with ``culture conversion'' arguments that point to the elective affinity between popular culture and weak-tie formation and maintenance processes \citep{lizardo06a}.  Activities associated with the higher prestige fine and performance arts, on the other hand, appear to be associated with having been contacted directly by the employer. While both of these can be considered prototypical ``weak-tie'' activation episodes, they differ in the underlying phenomenology and mechanism:  the first is to be expected of individuals who have networks rich in fleeting acquaintances directed by themselves towards others (outdegree) while the latter is characteristic of individuals who are likely to be known or sought out by a large number of potential contacts (indegree). Omnivores---who consume both popular and high status culture---may thus be able to enjoy \emph{both} of these network advantages when exploiting social resources for instrumental purposes. In this manner cultural variety affords the person with resource complementarity as manifested in both lots of outgoing ties directed to the person \emph{toward} diverse others and in lots of connections coming to the individual \emph{from} others, most importantly those located in positions of influence or authority \citep{yakubovich05, bian97, podolny_baron97}.

This analysis contributes to our understanding of the role of culture in processes of stratification by connecting culture consumption variety not only to the static \emph{possession} of virtual resources in the form of larger and more varied networks of social relations, but by providing empirical evidence of a link between culture consumption variety and the likelihood of attaining desirable outcomes through the \emph{activation} of personal relationships.  This is important because it is through the activation of resources embedded in social relations that social capital is theorized as going from a potential to a real resource \citep{lin99, lin01a}. These results are consistent with the general proposition that cultural taste variety ``counts'' as cultural capital in contemporary social settings precisely due to the intimate association of culture consumption and sociability.  Heterogeneity in the opportunities and likelihood to engage in this type of sociable interaction, can then be traced differences in rates of engagement in culture consumption \citep{witte_ryan04}.  This points to the building of social relationships through face-to-interaction as mediated by cultural content as one of the key mechanisms linking cultural stratification and status attainment outcomes.
